\section{Related Work}
\subsection{Commercial AI agents and Vertical Workflows}
Recent advances in foundation models have enabled AI agents to gradually acquire planning, tool-use, file-processing, and multi-step execution capabilities \citep{yao2023react, schick2023toolformer}, pushing AI systems from general conversational assistants toward workflow-oriented products. In practice, commercial AI products increasingly package these capabilities into vertical workflows such as document processing, data analysis, financial research, enterprise operations, and knowledge work. This trend suggests that real-world AI value is not determined only by isolated model capabilities, but also by whether these capabilities can be productized into useful work products. However, existing academic benchmarks rarely use commercial adoption itself as a source of task discovery. StartupBench addresses this gap by deriving tasks from market-validated AI-native startups and their deep users, moving task construction from researcher-defined settings toward demand-driven workflow construction.

\subsection{End-to-end evaluation of AI agents}
Agent benchmarks have evolved from capability evaluation toward end-to-end task completion. Early benchmarks mainly focus on tool use, multi-step reasoning, and general assistant abilities, as in AgentBench and GAIA \citep{liu2024agentbench, mialon2024gaia}. Later work extends evaluation to more realistic interactive environments, including web interaction and computer-use tasks \citep{zhou2024webarena, xieosworld}. More recent benchmarks further move toward professional work and economically meaningful tasks, covering enterprise software, simulated workplaces, occupational tasks, and long-horizon workflows \citep{drouin2403workarena, xu2026theagentcompany, patwardhan2025gdpval, hu2026occubenchevaluatingaiagents, zhu2026workflowgymlonghorizonevaluationcomputeruse,meng2026dvworld}. These works reflect a clear transition from isolated capability tests to realistic task-completion evaluation. StartupBench follows this direction, but differs in its task source: rather than relying primarily on public environments, occupational taxonomies, or simulated workplaces, it derives tasks from AI-native startup products, ToB user needs, and expected business deliverables.
\subsection{Agent-as-a-Judge for complex deliverable evaluation}
As agent tasks shift from short answers to heterogeneous deliverables such as reports, spreadsheets, slides, code patches, and structured files, evaluation must move from capability-centric scoring to deliverable-centric assessment\cite{you2026agent}. LLM-as-a-Judge has been widely used for open-ended response evaluation, with MT-Bench and Prometheus studying preference-based or rubric-conditioned judging \citep{zheng2023judging, kim2024prometheus}. For agentic workflows, however, judging often requires evidence retrieval, file inspection, state verification, and constraint checking over final artifacts. Agent-as-a-Judge extends LLM judging by equipping evaluators with tools and environment interaction, while AJ-Bench further studies judge agents on information acquisition, state verification, and process verification \citep{zhuge2024agent, shi2026aj}. StartupBench adopts this deliverable-centric evaluation paradigm by combining expert rubrics with automated judge agents that inspect original artifacts and evidence views, evaluating whether agents can transform model capabilities into usable work products.


%\subsection{Hello World}